\documentclass[10pt]{article}
\usepackage[french]{babel}
\usepackage[utf8]{inputenc}
\usepackage[T1]{fontenc}
\usepackage{amsmath}
\usepackage{amsfonts}
\usepackage{amssymb}
\usepackage[version=4]{mhchem}
\usepackage{stmaryrd}

\begin{document}
\section*{Analyse Numérique - correction TD n \({ }^{\boldsymbol{\circ}} \mathbf{4}\) 20-03-2025}
\section*{Exercice 1:}
\begin{enumerate}
  \item Le problème sera bien posé si la condition initiale vérifie bien les conditions aux limites qui ici, sont toutes les deux, des conditions de Neumann.\\
\(\theta_{0}=\alpha x^{2}+1 \quad \frac{\partial \theta_{0}}{\partial x}=2 \alpha x\)\\
Pour \(x=0 \quad \frac{\partial \theta_{0}}{\partial x}=0\)\\
Pour \(x=L \quad \frac{\partial \theta_{0}}{\partial x}=2 \alpha L\)\\
Dans le cas présent, les conditions initiales vérifient bien les conditions aux limites. Donc le problème est bien posé.
  \item Résolvons numériquement\\
\(\left.\frac{\partial \theta}{\partial t}=K \frac{\partial^{2} \theta}{\partial x^{2}}, \quad x \in\right] 0, L[, \forall t>0\)\\
Le terme de droite de l'équation est linéaire. Appelons le L , avec \(L=K \frac{\partial^{2} \theta}{\partial x^{2}}\)
\end{enumerate}

\section*{Schéma en temps}
L'équation étant linéaire, la résolution temporelle peut se faire avec un schéma implicite.\\
Choisissons le schéma implicite d'Adams qui s'écrit :\\
\(\theta_{i}^{n+1}=\theta_{i}^{n}+\Delta t\left(\frac{3}{4} L_{i}^{n+1}+\frac{1}{4} L_{i}^{n-1}\right)\)

\section*{Remarque}
Cette formule récursive pose un problème pour \(n=0\). En effet :\\
\(\theta_{i}^{1}=\theta_{i}^{0}+\Delta t\left(\frac{3}{4} L_{i}^{1}+\frac{1}{4} L_{i}^{-1}\right)\), le terme \(\frac{1}{4} L_{i}^{-1} \mathrm{n}^{\prime}\) a pas de sens. Il faut donc utiliser un autre schéma pour obtenir \(\theta_{i}^{1}\). Ce sera le schéma d'Euler Implicite \(\theta_{i}^{1}=\theta_{i}^{0}+\Delta t L_{i}^{1}\)

D'où le schéma temporel général :

\[
\begin{aligned}
t^{1} \quad(n=0) \quad & \theta_{i}^{1}=\theta_{i}^{0}+\Delta t L_{i}^{1} \\
& \theta_{i}^{1}-\Delta t L_{i}^{1}=\theta_{i}^{0}
\end{aligned}
\]

Pour \(n>0 \quad \theta_{i}^{n+1}=\theta_{i}^{n}+\Delta t\left(\frac{3}{4} L_{i}^{n+1}+\frac{1}{4} L_{i}^{n-1}\right)\)

\[
\begin{array}{ll} 
& \theta_{i}^{n+1}-\frac{3 \Delta t}{4} L_{i}^{n+1}=\theta_{i}^{n}+\frac{\Delta t}{4} L_{i}^{n-1} \\
t^{2} \quad(n=1) \quad \theta_{i}^{2}-\frac{3 \Delta t}{4} L_{i}^{2}=\theta_{i}^{1}+\frac{\Delta t}{4} L_{i}^{0} \\
t^{3} \quad(n=2) \quad \theta_{i}^{3}-\frac{3 \Delta t}{4} L_{i}^{3}=\theta_{i}^{2}+\frac{\Delta t}{4} L_{i}^{1}
\end{array}
\]

\section*{Schéma en espace}
Pour le schéma en espace, nous choisissons des schémas aux différence finies.

\[
\frac{\partial^{2} \theta}{\partial x^{2}} \approx \frac{\theta_{i+1}^{n}-2 \theta_{i}^{n}+\theta_{i-1}^{n}}{\Delta x^{2}}, \quad i=2, \ldots, N-1
\]

Avec ce type de schéma, les conditions aux limites, s'écrivent :

\[
\begin{aligned}
& i=1,\left.\quad \frac{\partial \theta}{\partial x}\right|_{x=0} \approx \frac{\theta_{2}^{n}-\theta_{1}^{n}}{\Delta x}=0, \quad \forall n \\
& i=N,\left.\quad \frac{\partial \theta}{\partial x}\right|_{x=L} \approx \frac{\theta_{N}^{n}-\theta_{N-1}^{n}}{\Delta x}=2 \alpha L, \quad \forall n
\end{aligned}
\]

\begin{itemize}
  \item Pour la première itération \((n=0)\), on a le problème discrétisé suivant :
\end{itemize}

\[
\begin{aligned}
& \theta_{i}^{1}-K \Delta t\left(\frac{\theta_{i+1}^{1}-2 \theta_{i}^{1}+\theta_{i-1}^{1}}{\Delta x^{2}}\right)=\theta_{i}^{0} \\
& \text { Posons } \beta=\frac{K \Delta t}{\Delta x^{2}}, \quad i=2, \ldots, N-1 \\
& \theta_{i}^{1}+2 \beta \theta_{i}^{1}-\beta \theta_{i+1}^{1}-\beta \theta_{i-1}^{1}=\theta_{i}^{0} \\
& -\beta \theta_{i+1}^{1}+(1+2 \beta) \theta_{i}^{1}-\beta \theta_{i-1}^{1}=\theta_{i}^{0} \\
& i=1, \quad \frac{\theta_{2}^{1}-\theta_{1}^{1}}{\Delta x}=0 \\
& i=N, \quad \frac{\theta_{N}^{1}-\theta_{N-1}^{1}}{\Delta x}=2 \alpha L
\end{aligned}
\]

\begin{itemize}
  \item Pour \(\mathrm{n}>0\), on a le problème discrétisé suivant :
\end{itemize}

\[
\begin{aligned}
& \theta_{i}^{n+1}-K \frac{3 \Delta t}{4}\left(\frac{\theta_{i+1}^{n+1}-2 \theta_{i}^{n+1}+\theta_{i-1}^{n+1}}{\Delta x^{2}}\right)=\theta_{i}^{n}+K \frac{\Delta t}{4}\left(\frac{\theta_{i+1}^{n-1}-2 \theta_{i}^{n-1}+\theta_{i-1}^{n-1}}{\Delta x^{2}}\right) \\
& \text { Posons } \sigma=\frac{3 K \Delta t}{4 \Delta x^{2}}, \quad i=2, \ldots, N-1 \\
& \theta_{i}^{n+1}+2 \sigma \theta_{i}^{n+1}-\sigma \theta_{i+1}^{n+1}-\sigma \theta_{i-1}^{n+1}=S_{i}^{n} \\
& -\sigma \theta_{i+1}^{n+1}+(1+2 \sigma) \theta_{i}^{n+1}-\sigma \theta_{i-1}^{n+1}=S_{i}^{n} \quad \text { avec } \\
& \mathrm{S}_{\mathrm{i}}^{\mathrm{n}}=\theta_{i}^{n}+\frac{\beta}{4} \theta_{i+1}^{n-1}-\frac{\beta}{2} \theta_{i}^{n-1}+\frac{\beta}{4} \theta_{i-1}^{n-1} \\
& \text { et } \quad \frac{\theta_{2}^{n+1}-\theta_{1}^{n+1}}{\Delta x}=0 \quad \text { et } \quad \frac{\theta_{N}^{n+1}-\theta_{N-1}^{n+1}}{\Delta x}=2 \alpha L
\end{aligned}
\]

\begin{enumerate}
  \setcounter{enumi}{2}
  \item Ecrivons le problème matriciel pour 6 points de discrétisation, y compris les bords.
\end{enumerate}

Pour calculer \(\theta_{i}^{1}, \quad i=1, \ldots, 6\)

\[
\left[\begin{array}{cccccc}
-\frac{1}{\Delta x} & \frac{1}{\Delta x} & 0 & 0 & 0 & 0 \\
-\beta & (1+2 \beta) & -\beta & 0 & 0 & 0 \\
0 & -\beta & (1+2 \beta) & -\beta & 0 & 0 \\
0 & 0 & -\beta & (1+2 \beta) & -\beta & 0 \\
0 & 0 & 0 & -\beta & (1+2 \beta) & -\beta \\
0 & 0 & 0 & 0 & -\frac{1}{\Delta x} & \frac{1}{\Delta x}
\end{array}\right]\left[\begin{array}{l}
\theta_{1}^{1} \\
\theta_{2}^{1} \\
\theta_{3}^{1} \\
\theta_{4}^{1} \\
\theta_{5}^{1} \\
\theta_{6}^{1}
\end{array}\right]=\left[\begin{array}{c}
0 \\
\theta_{2}^{0} \\
\theta_{3}^{0} \\
\theta_{4}^{0} \\
\theta_{5}^{0} \\
2 \alpha L
\end{array}\right]
\]

Pour calculer \(\theta_{i}^{n+1}, \quad i=1, \ldots, 6\)

\[
\left[\begin{array}{cccccc}
-\frac{1}{\Delta x} & \frac{1}{\Delta x} & 0 & 0 & 0 & 0 \\
-\sigma & (1+2 \sigma) & -\sigma & 0 & 0 & 0 \\
0 & -\sigma & (1+2 \sigma) & -\sigma & 0 & 0 \\
0 & 0 & -\sigma & (1+2 \sigma) & -\sigma & 0 \\
0 & 0 & 0 & -\sigma & (1+2 \sigma) & -\sigma \\
0 & 0 & 0 & 0 & -\frac{1}{\Delta x} & \frac{1}{\Delta x}
\end{array}\right]\left[\begin{array}{c}
\theta_{1}^{n+1} \\
\theta_{2}^{n+1} \\
\theta_{3}^{n+1} \\
\theta_{4}^{n+1} \\
\theta_{5}^{n+1} \\
\theta_{6}^{n+1}
\end{array}\right]=\left[\begin{array}{c}
0 \\
S_{2}^{n} \\
S_{3}^{n} \\
S_{4}^{n} \\
S_{5}^{n} \\
2 \alpha L
\end{array}\right]
\]

\begin{enumerate}
  \setcounter{enumi}{3}
  \item Pour résoudre le problème par une méthode itérative (Gauss-Seidel par exemple), il faut que la matrice soit à diagonale strictement dominante :
\end{enumerate}

\[
\begin{aligned}
\sum_{j=1}^{N}\left|a_{i j}\right|<\left|a_{i i}\right| \quad \text { avec } \mathrm{i} & =1, \ldots, \mathrm{~N} \\
\mathrm{j} & =1, \ldots, \mathrm{~N} \text { et } \mathrm{j} \neq \mathrm{i}
\end{aligned}
\]

Mis à part la première ligne et la dernière ligne de la matrice pour lesquelles :\\
\(\sum_{j=1}^{N}\left|a_{1 j}\right|=\left|a_{11}\right|=\frac{1}{\Delta x} \quad\) et\\
\(\sum_{j=1}^{N}\left|a_{N j}\right|=\left|a_{N N}\right|=\frac{1}{\Delta x}\)\\
Ses termes vérifient bien \(\sum_{j=1}^{N}\left|a_{i j}\right|=2 \sigma<\left|a_{i i}\right|=1+2 \sigma\)\\
5. On vérifie tout d'abord que ces nouvelles conditions sont bien posées : conditions de Neumann.\\
\(\theta_{0}=3 x^{2}+1 \quad \frac{\partial \theta_{0}}{\partial x}=6 x\)\\
Pour \(x=0 \quad \theta_{0}=1\) (condition de Dirichlet)\\
Pour \(x=L \quad \frac{\partial \theta_{0}}{\partial x}=6 L\) (condition de Neumann)\\
Les conditions initiales vérifient bien les conditions aux limites, donc le problème est de nouveau bien posé.

Les équations matricielles seront modifiées de la sorte :

Pour calculer \(\theta_{i}^{1}, \quad i=1, \ldots, 6\)

\[
\left[\begin{array}{cccccc}
1 & 0 & 0 & 0 & 0 & 0 \\
-\beta & (1+2 \beta) & -\beta & 0 & 0 & 0 \\
0 & -\beta & (1+2 \beta) & -\beta & 0 & 0 \\
0 & 0 & -\beta & (1+2 \beta) & -\beta & 0 \\
0 & 0 & 0 & -\beta & (1+2 \beta) & -\beta \\
0 & 0 & 0 & 0 & -\frac{1}{\Delta x} & \frac{1}{\Delta x}
\end{array}\right]\left[\begin{array}{c}
\theta_{1}^{1} \\
\theta_{2}^{1} \\
\theta_{3}^{1} \\
\theta_{4}^{1} \\
\theta_{5}^{1} \\
\theta_{6}^{1}
\end{array}\right]=\left[\begin{array}{c}
1 \\
\theta_{2}^{0} \\
\theta_{3}^{0} \\
\theta_{4}^{0} \\
\theta_{5}^{0} \\
6 L
\end{array}\right]
\]

Pour calculer \(\theta_{i}^{n+1}, \quad i=1, \ldots, 6\)

\[
\left[\begin{array}{cccccc}
1 & 0 & 0 & 0 & 0 & 0 \\
-\sigma & (1+2 \sigma) & -\sigma & 0 & 0 & 0 \\
0 & -\sigma & (1+2 \sigma) & -\sigma & 0 & 0 \\
0 & 0 & -\sigma & (1+2 \sigma) & -\sigma & 0 \\
0 & 0 & 0 & -\sigma & (1+2 \sigma) & -\sigma \\
0 & 0 & 0 & 0 & -\frac{1}{\Delta x} & \frac{1}{\Delta x}
\end{array}\right]\left[\begin{array}{c}
\theta_{1}^{n+1} \\
\theta_{2}^{n+1} \\
\theta_{3}^{n+1} \\
\theta_{4}^{n+1} \\
\theta_{5}^{n+1} \\
\theta_{6}^{n+1}
\end{array}\right]=\left[\begin{array}{c}
1 \\
S_{2}^{n} \\
S_{3}^{n} \\
S_{4}^{n} \\
S_{5}^{n} \\
6 L
\end{array}\right]
\]

\section*{Exercice 2:}
L'équation de Burgers s'écrit :

\[
\left.\frac{\partial u}{\partial t}=-u \frac{\partial u}{\partial x}+v \frac{\partial^{2} u}{\partial x^{2}}=H+L \quad x \in\right] 0, L[
\]

avec

\[
H=-u \frac{\partial u}{\partial x} \text { et } L=v \frac{\partial^{2} u}{\partial x^{2}}
\]

\begin{enumerate}
  \item En appliquant un schéma d'Euler explicite pour la partie non linéaire (H) et un schéma d'Euler implicite pour la partie linéaire (L), on obtient
\end{enumerate}

\[
\frac{u_{i}^{n+1}-u_{i}^{n}}{\Delta t}=H_{i}^{n}+L_{i}^{n+1}=\left(-u \frac{\partial u}{\partial x}\right)_{i}^{n}+\left(v \frac{\partial^{2} u}{\partial x^{2}}\right)_{i}^{n+1}, \quad 2 \leq i \leq N-1
\]

En prenant des schémas des différences-finies centrées d'ordre 2 pour la dérivée première et la dérivée seconde, on a :\\
\(\left(\frac{\partial u}{\partial x}\right)_{i}^{n}=\left(\frac{u_{i+1}^{n}-u_{i-1}^{n}}{2 \Delta x}\right)\), on a alors

\[
H_{i}^{n}=\left(-u \frac{\partial u}{\partial x}\right)_{i}^{n}=-u_{i}^{n}\left(\frac{u_{i+1}^{n}-u_{i-1}^{n}}{2 \Delta x}\right)
\]

\[
\begin{aligned}
& \left(\frac{\partial^{2} u}{\partial x^{2}}\right)_{i}^{n+1}=\left(\frac{u_{i+1}^{n+1}-2 u_{i}^{n+1}+u_{i-1}^{n+1}}{\Delta x^{2}}\right), \text { on a } \\
& L_{i}^{n+1}=\left(v \frac{\partial^{2} u}{\partial x^{2}}\right)_{i}^{n+1}=v\left(\frac{u_{i+1}^{n+1}-2 u_{i}^{n+1}+u_{i-1}^{n+1}}{\Delta x^{2}}\right)
\end{aligned}
\]

En remplaçant dans le schéma, on obtient

\[
\begin{aligned}
& \frac{u_{i}^{n+1}-u_{i}^{n}}{\Delta t}=-u_{i}^{n}\left(\frac{u_{i+1}^{n}-u_{i-1}^{n}}{2 \Delta x}\right)+v\left(\frac{u_{i+1}^{n+1}-2 u_{i}^{n+1}+u_{i-1}^{n+1}}{\Delta x^{2}}\right), \quad 2 \leq i \leq N-1 \\
& u_{i}^{n+1}-v \Delta t\left(\frac{u_{i+1}^{n+1}-2 u_{i}^{n+1}+u_{i-1}^{n+1}}{\Delta x^{2}}\right)=u_{i}^{n}-\Delta t u_{i}^{n}\left(\frac{u_{i+1}^{n}-u_{i-1}^{n}}{2 \Delta x}\right), \quad 2 \leq i \leq N-1
\end{aligned}
\]

Si on pose

\[
\begin{aligned}
& \sigma=\frac{v \Delta t}{\Delta x^{2}} \\
& -\sigma u_{i+1}^{n+1}+(1+2 \sigma) u_{i}^{n+1}-\sigma u_{i-1}^{n+1}=u_{i}^{n}-\Delta t u_{i}^{n}\left(\frac{u_{i+1}^{n}-u_{i-1}^{n}}{2 \Delta x}\right), \quad 2 \leq i \leq N-1 \\
& -\sigma u_{i+1}^{n+1}+(1+2 \sigma) u_{i}^{n+1}-\sigma u_{i-1}^{n+1}=S_{i}^{n}, \quad 2 \leq i \leq N-1
\end{aligned}
\]

avec

\[
S_{i}^{n}=u_{i}^{n}-\Delta t u_{i}^{n}\left(\frac{u_{i+1}^{n}-u_{i-1}^{n}}{2 \Delta x}\right)
\]

\begin{enumerate}
  \setcounter{enumi}{1}
  \item A quelle condition le problème est-il bien posé ? On choisit des conditions aux limites de Dirichlet \(\mathrm{u}=\mathrm{u} 0=\) constante.
\end{enumerate}

Le problème est bien posé si et seulement si les conditions initiales vérifient les conditions aux limites. Les conditions aux limites choisies ici sont des conditions de Dirichlet et sont constantes, on a alors :

\[
u_{1}^{n+1}=u_{0} \text { et } u_{N}^{n+1}=u_{0}, \quad \forall n
\]

Afin de vérifier les conditions aux limites, il suffit de prendre des conditions initiales égales aux conditions aux limites

\[
u_{1}^{0}=u_{N}^{0}=u_{0}
\]

Pour les autres termes \(u_{i}^{0}, \quad 2 \leq i \leq N-1\), il faut définir un état initial.\\
3. Ecrire le problème matriciel ainsi obtenu pour une discrétisation spatiale de 6 points (y compris les bornes du domaine de définition) régulièrement espacés. Suggérer une méthode de résolution et justifier la.

Pour six points d'espace ( \(\mathrm{N}=6\) ), on a le système suivant

\[
\left[\begin{array}{cccccc}
1 & 0 & 0 & 0 & 0 & 0 \\
-\sigma & (1+2 \sigma) & -\sigma & 0 & 0 & 0 \\
0 & -\sigma & (1+2 \sigma) & -\sigma & 0 & 0 \\
0 & 0 & -\sigma & (1+2 \sigma) & -\sigma & 0 \\
0 & 0 & 0 & -\sigma & (1+2 \sigma) & -\sigma \\
0 & 0 & 0 & 0 & 0 & 1
\end{array}\right]\left[\begin{array}{c}
u_{1}^{n+1} \\
u_{2}^{n+1} \\
u_{3}^{n+1} \\
u_{4}^{n+1} \\
u_{5}^{n+1} \\
u_{6}^{n+1}
\end{array}\right]=\left[\begin{array}{c}
u_{0} \\
S_{2}^{n} \\
S_{3}^{n} \\
S_{4}^{n} \\
S_{5}^{n} \\
u_{0}
\end{array}\right]
\]

Le système matriciel peut être résolu par une méthode itérative (Jacobi ou Gauss-Seidel) si la matrice A est diagonale strictement dominante (condition suffisante mais pas nécessaire).

La matrice A est à diagonale strictement dominante si :

\[
\sum_{j=1}^{n}\left|a_{i j}\right|<\left|a_{i i}\right|, j \neq i, i=1, \ldots, n
\]

On a

\begin{itemize}
  \item pour la première ligne et dernière ligne \(1>0\) OK pour \(\mathrm{i}=2, \ldots, 5\)
  \item \(|1+2 \sigma|>|-\sigma|+|-\sigma| \Rightarrow(1+2 \sigma)>2 \sigma \quad\) ок
\end{itemize}

Donc le système matriciel peut être résolu par une méthode itérative


\end{document}